\chapter[Constante de Planck y escala interna de acción]{Constante de Planck y recta de acción: derivación determinantal de la escala \texorpdfstring{\(10^{-34}\)}{decádica -34}}
\label{chap:escala-accion}
\index{Planck, constante de}\index{acción!escala de Planck}

La construcción de la escala de acción comienza en un invariante entero. La
forma normal de Smith del bloque dodecafásico local produce un conúcleo de
orden \(16\); su norma sobre la sección constante \(\alphaHMT\), seguida por
una única actuación de \(\varphi\), sitúa la salida
\(\Sigma_H=\varphi\alphaHMT^{16}\) entre \(10^{-34}\) y \(10^{-33}\). De
este modo, la potencia decimal característica de la acción se obtiene antes
de elegir joules y segundos. La rama angular aporta después las dos
coordenadas orientadas ---previa al retorno y posterior al retorno---, y la
carta metrológica las expresa finalmente como \(h\) y \(\hbar\).

\begin{exactbox}[title={Resultado rector: del conúcleo local a la escala de acción}]
La forma de Smith de \(H_4\), el orden de su conúcleo y las desigualdades que
determinan la década son exactos. El resultado completo es
\textsc{demostrado con estructura de partida explícita}: se declaran el
bloque local, la sección invariante, la norma multiplicativa, una única
actuación de \(\varphi\) y el lector decimal de
\cref{chap:orden-determinantal-accion}. Las ablaciones identifican la función
de cada regla y distinguen el conúcleo local del conúcleo global de orden
\(4096\).
\end{exactbox}

\section{Dato heredado y proposición de transferencia}

Para una de las tres órbitas dodecafásicas paralelas, las Partes~I--III obtienen
\[
\operatorname{SNF}(H_4)=\operatorname{diag}(1,2,2,4),
\qquad
G_H:=\operatorname{coker}(H_4)
\cong\ZZ/2\ZZ\oplus\ZZ/2\ZZ\oplus\ZZ/4\ZZ,
\]
y, por tanto, \(|G_H|=16\). La localidad pertenece al tipo del lector: el
conúcleo de la transformación dodecafásica completa tiene orden
\(16^3=4096\), pero registra simultáneamente tres copias directas y no
sustituye al índice de una órbita.

Sea \(s_\alpha(g)=\alphaHMT\) la sección constante sobre \(G_H\). Con la
norma local
\[
\mathcal N_{G_H}(s):=\prod_{g\in G_H}s(g)
\]
y una sola actuación de \(\varphi\) sobre la base, se recupera, con una
etiqueta local para facilitar las referencias de la Parte~V,
\begin{equation}
\label{eq:sigma-accion}
\Sigma_H
:=\varphi\,\mathcal N_{G_H}(s_\alpha)
=\varphi\alphaHMT^{16}.
\end{equation}
Esta es la misma sección definida en \cref{eq:sigma-h-local}, no una segunda
construcción.

\begin{theorem}[Transferencia del orden determinantal]
\label{thm:orden-accion}
Fijadas las primitivas del lector determinantal local, la sección de
\cref{eq:sigma-accion} satisface
\[
\alphaHMT^{16}<10^{-34}
<\varphi\alphaHMT^{16}<10^{-33},
\qquad
\varepsilon_H:=\left\lfloor\log_{10}\Sigma_H\right\rfloor=-34.
\]
Si \(\mathcal L_S\) es una recta real de acción con base
\(\mathcal U_S\), la asignación
\[
\tau_S:\RR_{>0}\longrightarrow\mathcal L_S,
\qquad x\longmapsto x\,\mathcal U_S,
\]
transporta ese orden sin modificarlo. En particular, el exponente decimal
de las coordenadas de acción definidas más abajo es \(-34\).
\end{theorem}

\begin{proof}
Las desigualdades exactas y el valor del orden son el contenido de
\cref{thm:exponente-accion-local}; allí se reducen a comparaciones enteras a
partir del decimal racional construido de \(\alphaHMT\). La aplicación
\(\tau_S\) cambia el codominio —de un escalar positivo a una recta de
magnitudes con base declarada—, pero conserva la coordenada escalar. Por
consiguiente, no altera su orden decimal.
\end{proof}

La evaluación heredada es
\[
\Sigma_H
=1.046243838104756580184229208303089\ldots\times10^{-34},
\]
y el número normalizado
\[
B_H:=10^{34}\Sigma_H
=1.046243838104756580184229208303089\ldots
\]
pertenece a \([1,10)\). El factor \(10^{34}\) se reconstruye a partir del
orden; no se introduce como umbral o normalización previa.

\section{Necesidad del conúcleo local}

Las cuatro ablaciones exactas son
\[
\begin{array}{c|c}
\text{regla evaluada}&\text{exponente decádico}\\
\hline
\alphaHMT^{16}\ \text{sin }\varphi&-35\\
\varphi\alphaHMT^{15}&-32\\
\varphi\alphaHMT^{17}&-37\\
\varphi\alphaHMT^{4096}&-8753.
\end{array}
\]
Así, \(|G_H|=16\) fija el exponente de \(\alphaHMT\) una vez escogido el
lector local, y las desigualdades fijan entonces el orden \(-34\). Estos
controles excluyen la omisión de \(\varphi\), los exponentes vecinos y la
sustitución del conúcleo local por el global. No constituyen una
clasificación de todos los lectores admisibles ni una prueba de unicidad
previa de la sección constante, de la actuación única de \(\varphi\) o de la
carta decimal.

\section{Acoplamiento posterior con las coordenadas de acción}

Escribamos
\[
\varepsilon_H:=\left\lfloor\log_{10}\Sigma_H\right\rfloor=-34.
\]
Una vez construidos por la rama angular el carácter previo al retorno \(H_5\)
y la coordenada reconstruida
\[
\etaRet=\frac{\Cstar}{2}-\alphaAngle,
\]
ambos reciben el mismo orden en \(\mathcal L_S\):
\[
\hbar_{\HMT}^{\rm pre}
:=H_5\,10^{\varepsilon_H}\mathcal U_S,
\qquad
\hbar_{\HMT}^{\rm ret}
:=\etaRet\,10^{\varepsilon_H}\mathcal U_S.
\]
Numéricamente,
\[
\hbar_{\HMT}^{\rm pre}
=1.054571817646154469625821829433059\ldots
\times10^{-34}\mathcal U_S,
\]
\[
\hbar_{\HMT}^{\rm ret}
=1.054571816915039061133690584724300\ldots
\times10^{-34}\mathcal U_S.
\]
El dato \(H_5\) es anterior al retorno y no se identifica con \(\etaRet\).
Tampoco interviene en la demostración del exponente \(-34\).

Las coordenadas no reducidas correspondientes son
\[
h_{\HMT}^{\rm pre}
=2\pi H_5\,10^{-34}\mathcal U_S
=6.626070149999987926001110349899474\ldots
\times10^{-34}\mathcal U_S,
\]
\[
h_{\HMT}^{\rm ret}
=2\pi\etaRet\,10^{-34}\mathcal U_S
=6.626070145406254333510749847592879\ldots
\times10^{-34}\mathcal U_S.
\]
La comparación
\[
2\pi H_5-6.62607015
=-1.2073998889650100526\ldots\times10^{-14}
\]
es un control externo posterior sobre la mantisa; no es una igualdad con el
valor definitorio del SI ni una entrada del lector.

\begin{center}
\begin{tabularx}{0.96\textwidth}{@{}L{0.22\textwidth}L{0.36\textwidth}Y@{}}
\toprule
Objeto & Coordenada obtenida & Función matemática y metrológica\\
\midrule
Escala determinantal &
\(\Sigma_H=1.0462438381047565\ldots\times10^{-34}\) &
Fija internamente el orden decimal de la recta de acción.\\
Acción reducida previa &
\(\hbar_{\rm HMT}^{\rm pre}=H_5\,10^{-34}\mathcal U_S\) &
Conserva la coordenada angular anterior al retorno.\\
Acción reducida posterior &
\(\hbar_{\rm HMT}^{\rm ret}=\etaRet\,10^{-34}\mathcal U_S\) &
Publica la coordenada post-retorno sobre la misma recta.\\
Acción de ciclo completo &
\(h_{\rm HMT}^{\rm pre}=6.6260701499999879\ldots\times10^{-34}\mathcal U_S\) &
Transporta la acción local por una vuelta de fase mediante \(h=2\pi\hbar\).\\
Carta SI &
\(h=6.62607015\times10^{-34}\ \mathrm{J\,s}\) &
Coordenada definitoria externa usada para el reconocimiento metrológico
\cite{BIPMSI}.\\
\bottomrule
\end{tabularx}
\end{center}

\section{\texorpdfstring{\(h\)}{h} y
\texorpdfstring{\(\hbar\)}{h barra}: acción por ciclo completo y por radián}
\label{sec:accion-vuelta-radian-rev6}
\index{acción!por vuelta}
\index{acción!por radián}

Las identidades
\[
 E=h\nu=\hbar\omega,\qquad
 \omega=2\pi\nu,\qquad
 h=2\pi\hbar
\]
no describen dos cuantos. Describen una misma acción en dos cartas. La
frecuencia \(\nu\) cuenta vueltas completas por unidad de tiempo; la
frecuencia angular \(\omega\) cuenta fase en radianes por unidad de tiempo.
En consecuencia, \(h\) mide la acción asociada al cierre de una vuelta y
\(\hbar\) mide la acción por unidad de fase. La reducción no es sólo una
abreviatura: transporta la holonomía global al generador local de su
propagación.

La misma distinción explica el uso coordinado de grados y radianes. El ciclo
angular completo de \(360^\circ\) hace visibles las particiones finitas
\[
 12\cdot30^\circ=4\cdot90^\circ
 =3\cdot120^\circ=360^\circ
\]
y los retornos de \(270^\circ\) y \(360^\circ\). El radián es, en cambio,
la coordenada aditiva en la que la exponencial y la derivada se escriben sin
insertar en cada operación una constante de conversión. Ambas cartas se
relacionan por
\[
 \mathbb R/\mathbb Z
 \xrightarrow{\;2\pi\;}
 \mathbb R/(2\pi\mathbb Z)
 \xrightarrow{\;180/\pi\;}
 \mathbb R/(360\mathbb Z).
\]
La composición conserva el punto de la vuelta: el grado publica su posición
en el calendario global y el radián linealiza su movimiento local.

Aplicado a las coordenadas construidas arriba,
\[
 h_{\rm HMT}^{\rm pre}
 =2\pi\,\hbar_{\rm HMT}^{\rm pre},\qquad
 h_{\rm HMT}^{\rm ret}
 =2\pi\,\hbar_{\rm HMT}^{\rm ret}.
\]
El factor \(2\pi\) no corrige retrospectivamente una mantisa. Recorre con la
acción local una vuelta completa. Por eso una amplitud
\(\exp(\mathrm iS/\hbar)\), una relación de conmutación y una cuantización
de holonomía comparan una acción con la unidad local de fase, mientras el
cierre global se expresa con \(h\).

\section{Dos ramas causales y su confluencia}

La rama angular y la rama determinantal parten de
\(\alphaHMT\), pero ninguna genera retrospectivamente a la otra. La primera
se descompone como
\[
\begin{gathered}
(\alphaHMT,\varphi;\,9,5,4,12,54,\text{ regla graduada})
\longrightarrow H_5,\\
\alphaHMT\longrightarrow A=1000\alphaHMT\,{}^\circ,
\qquad
(\alphaHMT,\pi)\longrightarrow D_A,\\
\mathcal F_\pi\longrightarrow\rho_\pi=169,\qquad
(\rho_\pi,D_A,\alphaHMT)\longrightarrow\mathcal C_\pi(\alphaHMT),\\
(\pi,\alphaHMT,H_5,\mathcal C_\pi(\alphaHMT))
\longrightarrow y_*\longrightarrow\Cstar,\\
(A,\Cstar)\longrightarrow q_\pm
\longrightarrow\bigl(S_{90}(q_\pm),S_{120}(q_\pm)\bigr)
\longrightarrow K_{6\to8}
\longrightarrow\Gamma_{6\to8}.
\end{gathered}
\]
La identificación posterior
\(\Gamma_{6\to8}\dashrightarrow\gamma_{\rm BI}\) pertenece a la interfaz
física. La rama determinantal es
\[
(H_4,\alphaHMT,\varphi;\text{ reglas del lector})
\longrightarrow G_H
\longrightarrow\Sigma_H
\longrightarrow\varepsilon_H=-34.
\]
Sólo al final, \(\varepsilon_H\) se combina con \(H_5\) y \(\etaRet\) para
dar coordenadas en \(\mathcal L_S\). Quedan excluidas tanto una flecha
\(A\to H_5\) como cualquier selección de \(C^*\) desde el SI o desde
\(\gamma_{\rm BI}\).

\section{Carta metrológica}

Elegir una unidad física es una operación posterior al teorema. Una carta
metrológica es un isomorfismo de rectas de magnitudes
\[
\iota_{\rm SI}:\mathcal L_S\longrightarrow\mathcal L_{\mathrm{J\,s}},
\qquad
\iota_{\rm SI}(\mathcal U_S)=\mathrm{J\,s}.
\]
La unidad pertenece al codominio de \(\iota_{\rm SI}\); el exponente
\(-34\) pertenece al teorema aritmético relativo. La yuxtaposición de un
símbolo de unidad no convierte por sí sola un número adimensional en una
magnitud física.

\begin{statusbox}[title={Jerarquía probatoria de la constante de Planck}]
El orden \(10^{-34}\) procede del determinante local y queda fijado antes de
la metrología. Las mantisas \(H_5\) y \(\etaRet\) proceden de la rama angular
y distinguen las orientaciones previa y posterior al retorno. La identidad
\(h=2\pi\hbar\) transporta ambas coordenadas entre acción por radián y acción
por ciclo. La aplicación \(\iota_{\rm SI}\) proporciona después la unidad
\(\mathrm{J\,s}\) y permite comparar la salida con la constante definitoria
del SI. Esta sucesión conserva en lugares distintos la derivación interna,
la realización dimensional y el contraste externo.
\end{statusbox}
